\documentclass[11pt,a4paper]{article}
\usepackage[T1]{fontenc}
\usepackage[utf8]{inputenc}
\usepackage[main=italian,provide=*]{babel}
\usepackage{amsmath,amssymb}
\usepackage{geometry}
\geometry{margin=2.5cm}
\usepackage{booktabs}
\usepackage{longtable}
\usepackage{seqsplit}
\usepackage{array}
\usepackage{hyperref}
\hypersetup{colorlinks=true,linkcolor=blue,citecolor=blue,urlcolor=blue}

\title{Trimero a catena aperta — fasi e documenti prodotti}
\author{Note di lavoro — Tesi triennale, Samuele Schiavi\\ Relatore: Prof.\ Alessandro Chiesa}
\date{}

\begin{document}
\maketitle

\section{Vincolo di partenza}

Su indicazione del relatore, la catena aperta $1$--$2$--$3$ include \textbf{solo}
le interazioni $1\leftrightarrow2$ e $2\leftrightarrow3$ (nessun legame
$1\leftrightarrow3$), con lo stesso percorso di lavoro già seguito per il
dimero: teoria esatta, VQE, fino a includere l'ansatz canonico. Interazione
DM e quantum simulation (Trotter, correlazioni) restano fuori scope finché
non esplicitamente richieste.

\section{Stato di avanzamento}

\begin{longtable}{@{}p{1.6cm}p{2.0cm}p{7.3cm}p{2.2cm}@{}}
\toprule
Fase & Stato & Documenti prodotti & \\
\midrule
\endhead

1 & \textbf{Completata} &
\texttt{\seqsplit{teoria\_trimero\_catena\_aperta.pdf}} (+\texttt{.tex}) --
teoria esatta: simmetria $P_{13}$, Casimir $S_{13}^2$, decomposizione di
Kambe, spettro chiuso a tre multipletti $A',B',C'$, campo critico
$b_c=3J$, incroci veri (non anticrossing), confronto strutturale con
l'anello.
&\\[4pt]
& &
\texttt{\seqsplit{animazione\_trimero\_catena\_aperta.html}} -- visualizzazione
interattiva dello spettro in campo (stessa grafica dell'anello, adattata:
un solo parametro $J$, geometria a catena con legame $1$--$3$ esplicitamente
assente).
&\\[10pt]

2 & \textbf{Completata} &
\texttt{\seqsplit{derivazione\_stati\_kambe\_trimero\_catena.pdf}} (+\texttt{.tex}) --
derivazione esplicita degli otto autostati nella base computazionale
$|q_1q_2q_3\rangle$, con gestione esplicita del riordino (la coppia di
classificazione $1,3$ non è adiacente nel registro fisico).
&\\[4pt]
& &
\texttt{\seqsplit{verifica\_stati\_catena.py}} -- verifica numerica indipendente
degli otto stati (autovettori esatti, ortonormalità, numeri quantici) a
precisione macchina.
&\\[10pt]

3 & \textbf{Completata} &
\texttt{\seqsplit{trimer\_chain\_exact.py}} -- benchmark esatto (Hamiltoniana
come \texttt{SparsePauliOp}, spettro analitico, stati di Kambe, proiettore
sul fondamentale, termine DM disponibile ma non usato), 8 self-test a
precisione macchina.
&\\[4pt]
& &
\texttt{\seqsplit{trimer\_chain\_exact.png}} -- spettro completo e
magnetizzazione del fondamentale, con distinzione esplicita fra il punto
convenzionale $b=0$ (fondamentale degenere) e l'incrocio fisico vero a
$b_c$.
&\\[10pt]

4 & \textbf{Completata} &
\texttt{\seqsplit{confronto\_ansatz\_entangler\_trimero\_catena.ipynb}} --
confronto sistematico di 10 ansatz VQE ($D=0$), ansatz canonico
individuato (\texttt{D: Ry-RBS-Ry}, 8 parametri, 4 gate a due qubit,
fidelity esatta ovunque), due limiti strutturali originali verificati
(PMA-1q non raggiunge $|111\rangle$ sulla catena pur riuscendoci
sull'anello; tetto esatto $5/6$ per la famiglia PMA-2q sotto $b_c$).
&\\[4pt]
& &
\texttt{\seqsplit{trimer\_chain\_ansatz\_sweep\_cache.json}},
\texttt{\seqsplit{confronto\_ansatz\_trimero\_catena\_fidelity.png}} -- cache
dei risultati e figura riassuntiva.
&\\[10pt]

5 & \textit{Non iniziata} &
Interazione DM sulla catena (unica combinazione compatibile con $P_{13}$:
$D_{12}=-D_{23}$) -- struttura già disponibile in
\texttt{trimer\_chain\_exact.py} (\texttt{dm\_term}, \texttt{trimer\_hamiltonian\_dm})
ma non utilizzata: non richiesta dal relatore allo stato attuale.
&\\[10pt]

-- & \textit{Non iniziata} &
Quantum simulation (Trotter) e correlazioni dinamiche sulla catena --
rimandate, per analogia diretta, a dopo il completamento delle stesse per
il trimero ad anello (vedi \texttt{convenzioni\_qubit\_progetto.pdf},
Parte 2, "Ancora da fare").
&\\

\bottomrule
\end{longtable}

\section{Documento di supporto trasversale}

\texttt{\seqsplit{convenzioni\_qubit\_progetto.pdf}} -- non specifico della catena,
ma aggiornato ad ogni fase: cataloga la convenzione sito$\leftrightarrow$qubit
di ciascun file del progetto (dimero, anello, catena), organizzato nello
stesso ordine cronologico della tesi, con il "perché" di ogni scelta.

\section{Prossimi passi}

In ordine di priorità naturale: (i) VQE con DM sulla catena, se richiesto
dal relatore; (ii) quantum simulation e correlazioni dinamiche
sull'anello (già in attesa, precedono cronologicamente la catena); (iii)
quantum simulation e correlazioni sulla catena, a seguire.

\end{document}
